\section{Largest Odd Number in a String}
\label{sec:str-largest-odd}

\problemstatement{Given a string of digits, return the largest odd number
that is a non-empty prefix substring (a prefix of the string), or the empty
string if none exists.}

\begin{patternbox}
A number is odd iff its \emph{last} digit is odd. Since we may only take a
prefix, the largest odd prefix ends at the \textbf{rightmost odd digit} --
so scan from the right for the first odd digit and return everything up to
it.
\end{patternbox}

\subsection*{Find the rightmost odd digit}

Any prefix of the string is a candidate number. To be odd it must end in an
odd digit, and to be largest it should be as long as possible. So scan from
the right end leftward for the first odd digit; the prefix ending there is
the answer. If no digit is odd, no odd prefix exists and the answer is
empty.

\begin{ideabox}[Oddness Depends Only on the Last Digit]
Because parity is decided by the final digit, the longest odd prefix is the
one ending at the rightmost odd digit. A single right-to-left scan finds it;
leading digits never change oddness, so length is maximised automatically.
\end{ideabox}

\subsection*{Algorithm}

\begin{enumerate}
  \item Scan from the last index leftward.
  \item At the first odd digit (index $i$), return the substring
        $[0,i]$.
  \item If none found, return the empty string.
\end{enumerate}

\subsection*{Dry run}

\dryrun{\texttt{"35427"} and \texttt{"5468"}.}

\begin{itemize}
  \item \texttt{"35427"}: last digit $7$ is odd $\to$ whole string
        \texttt{"35427"}.
  \item \texttt{"5468"}: rightmost odd is $5$ at index $0$ $\to$
        \texttt{"5"}.
\end{itemize}

\subsection*{C++ implementation}

\begin{lstlisting}
#include <bits/stdc++.h>
using namespace std;

string largestOddNumber(const string& s) {
    for (int i = (int)s.size() - 1; i >= 0; i--)
        if ((s[i] - '0') % 2 == 1) return s.substr(0, i + 1);   // odd prefix
    return "";
}

int main() {
    cout << "[" << largestOddNumber("35427") << "]\n";   // [35427]
    cout << "[" << largestOddNumber("5468") << "]\n";     // [5]
    return 0;
}
\end{lstlisting}

\begin{complexitybox}
\textbf{Time Complexity.} $O(n)$ -- one scan. \textbf{Space Complexity.}
$O(1)$ beyond the returned substring.
\end{complexitybox}

\begin{takeawaybox}
When a property depends only on a boundary character (here, oddness on the
last digit), scan toward that boundary rather than evaluating every prefix.
The rightmost-odd-digit observation collapses the problem to one pass.
\end{takeawaybox}
